\section{Setting, notation, and the rules of evidence}\label{sec:setting}

\subsection{The transport datum}

Throughout, $H$ is a finite-dimensional real or complex Hilbert space of dimension $n$,
$\ip{\cdot}{\cdot}$ its inner product, $\norm{\cdot}$ the operator norm and
$\norm{X}_F^2=\Tr(X^*X)$ the Frobenius (Hilbert--Schmidt) norm. $\PP$ always denotes an
\emph{orthogonal} projector, $\PP^*=\PP$, $\PP^2=\PP$, and $\QQ:=1-\PP$. $\Tr$ is the trace.
A \emph{transport datum} is a triple
\[
   \bigl(t\mapsto A(t),\ \PP(t),\ \Gap(t)\bigr),
\]
where $A(t)$ is a $\C^1$ family of self-adjoint operators, $\PP(t)$ is the orthogonal projector
onto a \emph{spectral cluster} $\shape(t)\subset\operatorname{spec}A(t)$, and $\Gap(t)>0$ is the
distance from $\shape(t)$ to the rest of the spectrum. The cluster is the \emph{admissible
subspace}; the gap $\Gap$ is the only quantity in the theory that is not derived from $A$ and
$\PP$ themselves.

\begin{remark}[Why a cluster and not an eigenvalue]\label{rem:cluster}
The theory is stated for clusters even where a single eigenvalue would do, because the two
behave differently under deformation (\S\ref{sec:hierarchy}). A single-eigenvalue cluster and a
rank-one projector carry the same information, and \S\ref{sec:hierarchy} shows that the
projector gains nothing; the gain appears only for clusters of two or more eigenvalues.
\end{remark}

\subsection{Invariants, and the one symbol per concept rule}

The corpus this paper belongs to has three different objects called $C$: a transport capacity,
a persistence overlap, and a boundary census. One symbol per concept, from here on.

\begin{center}
\small
\begin{tabularx}{\textwidth}{@{}l l >{\raggedright\arraybackslash}X@{}}
\toprule
symbol & object & defined in \\
\midrule
$\Omega$ & structural operator, a self-adjoint field & \S\ref{sec:redistribution} \\
$\Fr(\Omega,\PP)$ & redistribution operator $\PP\Omega\QQ+\QQ\Omega\PP$ & \S\ref{sec:redistribution} \\
$\pi(\PP_1,\PP_2)$ & persistence $\Tr(\PP_1\PP_2)$ & \S\ref{sec:setting} \\
$\rho$ & rigidity $\sigma_{\min}^{-2}$ & \S\ref{sec:shape} \\
$Q_{\mu\nu}$ & quantum geometric tensor & \S\ref{sec:geometry} \\
$g_{\mu\nu}$ & quantum metric $\operatorname{Re}Q_{\mu\nu}$ & \S\ref{sec:geometry} \\
$\omega_{\mu\nu}$ & curvature $-2\operatorname{Im}Q_{\mu\nu}$ & \S\ref{sec:geometry} \\
$\sigma_i$ & singular values of the transport differential & \S\ref{sec:shape} \\
$\kappa$ & cone slope & \S\ref{sec:shape} \\
\bottomrule
\end{tabularx}
\end{center}

The persistence of two projectors is the one invariant that needs no gap:
\begin{equation}\label{eq:persistence}
   \pi(\PP_1,\PP_2):=\Tr(\PP_1\PP_2)=\norm{\PP_2\PP_1}_F^2\ \in[0,\min(k_1,k_2)],
\end{equation}
which is \T\ and kernel-checked,\footnote{\lean{EmergentTransport/Basic.lean}:
\lean{persistence\_frobenius}, \lean{persistence\_nonneg}, and at rank one with value
$(u\cdot v)^2=\cos^2\theta$, \lean{persistence\_rank\_one}.} with the rank-one value
$(u\cdot v)^2=\cos^2\theta$; the principal-angle form $\sum_i\cos^2\theta_i$ for higher rank is
classical \Cc{} and is not used below.

\subsection{The five grades, and the weakest-link rule}

Every statement carries one of
\begin{center}
\begin{tabular}{@{}ll@{}}
\toprule
grade & meaning \\
\midrule
\T & proved in this paper (or kernel-checked, marked \lean{name}) \\
\V & verified by a named gate or script, with the file and the measured residual \\
\Cc & cited, with the hypotheses under which it holds stated \\
\Dd & data reported under a stated protocol, not measured here \\
\Hh & hypothesis or interpretation, not established \\
\bottomrule
\end{tabular}
\end{center}
Where a chain of reasoning mixes grades, the weakest link is named in the sentence that states
the conclusion. This is the corpus's own convention and it is the reason the paper can be
audited rather than believed.

\subsection{What is not assumed}

No metric on the space of parameters is assumed: the metric of \S\ref{sec:geometry} is
\emph{derived} from the projector family, and the statement that it is a metric is a theorem
(\S\ref{sec:geometry}), not a definition. No genericity is assumed except where stated
(simple eigenvalues, non-degeneracy of a gap). No smoothness beyond $\C^1$ is assumed except
where derivatives of higher order appear, and there it is stated. No physical system is
assumed; \S\ref{sec:scope} says exactly what is and is not claimed about the physical readings
the corpus attaches to the framework.
